% GENERATED FILE. Edit canonical lessons, then run npm run generate:books.
% canonical-source-hash: fnv1a64:9429d38df1cfdc90
% canonical-lessons: SA-C372-vighnah, SA-C372-badha, SA-C372-badhate, SA-C372-vilambayati, SA-C372-apanayati

\chapter{What Stands in the Way}
\label{ch:sa-what-stands-in-the-way}
\hlchaptermodality{372}

\begin{chapteropening}
\textbf{By the end of this chapter:} \emph{I can say what might get in the way and how to clear it.}

Complete the last lesson of chapter 372: I can say what might get in the way and how to clear it.
\end{chapteropening}

\section[\texorpdfstring{\sk{विघ्नः} (vighnaḥ)}{विघ्नः (vighnaḥ)}]{\sk{विघ्नः} (vighnaḥ) — an obstacle}

\label{lesson:SA-C372-vighnah}



\begin{quote}
Before the new one: say the Sanskrit for a time limit, then the Sanskrit for sets about.
\end{quote}

\subsection*{You'll want to know: \sk{विघ्नः}}
\textbf{\sk{विघ्नः}} — \emph{vighnaḥ} — \textquotedblleft{}an obstacle\textquotedblright{}.

\begin{quote}
\sk{यदि} \sk{विघ्नः} \sk{आगच्छति}, \sk{तर्हि} \sk{धैर्यं} \sk{न} \sk{त्यजतु।}
\end{quote}

\begin{quote}
\emph{yadi vighnaḥ āgacchati, tarhi dhairyaṃ na tyajatu.}
\end{quote}

\begin{quote}
If an obstacle comes, do not lose heart.
\end{quote}

A masculine noun.

\begin{cousinweb}
\textbf{\sk{विघ्नः}} (\emph{vighnaḥ}) is built on \emph{vi} and the root √han, \textquotedblleft{}to strike\textquotedblright{}; MW gives \emph{vi-ghna} \textquotedblleft{}an obstacle, an impediment, a hindrance\textquotedblright{} (MW).
\end{cousinweb}

\begin{grammarlens}[title={What you've built}]
One more word for saying what happens if.
\end{grammarlens}

\subsection*{Your turn}
\begin{itemize}
\raggedright
  \item \emph{Say it:} \emph{vighnaḥ}
  \item \emph{Say it:} \emph{vighnaḥ}, once more
  \item \emph{Say it:} \emph{yadi vighnaḥ āgacchati, tarhi dhairyaṃ na tyajatu.} — If an obstacle comes, do not lose heart.
  \item \emph{Recall:} say the Sanskrit for a time limit, then the Sanskrit for sets about, then say \emph{vighnaḥ} again
\end{itemize}

\begin{tcolorbox}[breakable,colback=teal!4,colframe=teal!35!black,title={Before you move on}]
What does \sk{विघ्नः} mean? (\textquotedblleft{}An obstacle\textquotedblright{}.) Say it once more.
\end{tcolorbox}
\section[\texorpdfstring{\sk{बाधा} (bādhā)}{बाधा (bādhā)}]{\sk{बाधा} (bādhā) — a hindrance; harm}

\label{lesson:SA-C372-badha}



\begin{quote}
Before the new one: say the Sanskrit for sets about, then the Sanskrit for an obstacle.
\end{quote}

\subsection*{You'll want to know: \sk{बाधा}}
\textbf{\sk{बाधा}} — \emph{bādhā} — \textquotedblleft{}a hindrance; harm\textquotedblright{}.

\begin{quote}
\sk{यदि} \sk{मार्गे} \sk{बाधा} \sk{अस्ति}, \sk{तर्हि} \sk{अन्यं} \sk{मार्गं} \sk{गच्छतु।}
\end{quote}

\begin{quote}
\emph{yadi mārge bādhā asti, tarhi anyaṃ mārgaṃ gacchatu.}
\end{quote}

\begin{quote}
If there is a hindrance on the road, take another road.
\end{quote}

A feminine noun.

\begin{cousinweb}
\textbf{\sk{बाधा}} (\emph{bādhā}) comes from the root √bādh, \textquotedblleft{}to press, to harass\textquotedblright{}; MW gives \emph{bādha} \textquotedblleft{}annoyance, obstacle, trouble\textquotedblright{} and the feminine \emph{bādhā} \textquotedblleft{}injury, hurt\textquotedblright{} (MW).
\end{cousinweb}

\begin{grammarlens}[title={What you've built}]
One more word for saying what happens if.
\end{grammarlens}

\subsection*{Your turn}
\begin{itemize}
\raggedright
  \item \emph{Say it:} \emph{bādhā}
  \item \emph{Say it:} \emph{bādhā}, once more
  \item \emph{Say it:} \emph{yadi mārge bādhā asti, tarhi anyaṃ mārgaṃ gacchatu.} — If there is a hindrance on the road, take another road.
  \item \emph{Recall:} say the Sanskrit for sets about, then the Sanskrit for an obstacle, then say \emph{bādhā} again
\end{itemize}

\begin{tcolorbox}[breakable,colback=teal!4,colframe=teal!35!black,title={Before you move on}]
What does \sk{बाधा} mean? (\textquotedblleft{}A hindrance; harm\textquotedblright{}.) Say it once more.
\end{tcolorbox}
\section[\texorpdfstring{\sk{बाधते} (bādhate)}{बाधते (bādhate)}]{\sk{बाधते} (bādhate) — hinders; troubles}

\label{lesson:SA-C372-badhate}



\begin{quote}
Before the new one: say the Sanskrit for an obstacle, then the Sanskrit for a hindrance.
\end{quote}

\subsection*{You'll want to know: \sk{बाधते}}
\textbf{\sk{बाधते}} — \emph{bādhate} — \textquotedblleft{}hinders; troubles\textquotedblright{}.

\begin{quote}
\sk{यदि} \sk{शब्दः} \sk{मां} \sk{बाधते}, \sk{तर्हि} \sk{अहम्} \sk{अन्यत्र} \sk{पठामि।}
\end{quote}

\begin{quote}
\emph{yadi śabdaḥ māṃ bādhate, tarhi aham anyatra paṭhāmi.}
\end{quote}

\begin{quote}
If the noise troubles me, I study elsewhere.
\end{quote}

Past (it …): \sk{शब्दः} \sk{मां} \sk{बाधितवान्}, \emph{śabdaḥ māṃ bādhitavān}.

\begin{cousinweb}
\textbf{\sk{बाधते}} (\emph{bādhate}) comes from the root √bādh, \textquotedblleft{}to press, to harass\textquotedblright{}; MW gives \emph{bādhita} \textquotedblleft{}pressed, oppressed\textquotedblright{} (MW).
\end{cousinweb}

\begin{grammarlens}[title={What you've built}]
One more word for saying what happens if.
\end{grammarlens}

\subsection*{Your turn}
\begin{itemize}
\raggedright
  \item \emph{Say it:} \emph{bādhate}
  \item \emph{Say it:} \emph{bādhate}, once more
  \item \emph{Say it:} \emph{yadi śabdaḥ māṃ bādhate, tarhi aham anyatra paṭhāmi.} — If the noise troubles me, I study elsewhere.
  \item \emph{Recall:} say the Sanskrit for an obstacle, then the Sanskrit for a hindrance, then say \emph{bādhate} again
\end{itemize}

\begin{tcolorbox}[breakable,colback=teal!4,colframe=teal!35!black,title={Before you move on}]
What does \sk{बाधते} mean? (\textquotedblleft{}Hinders; troubles\textquotedblright{}.) Say it once more.
\end{tcolorbox}
\section[\texorpdfstring{\sk{विलम्बयति} (vilambayati)}{विलम्बयति (vilambayati)}]{\sk{विलम्बयति} (vilambayati) — delays (something); holds up}

\label{lesson:SA-C372-vilambayati}



\begin{quote}
Before the new one: say the Sanskrit for a hindrance, then the Sanskrit for hinders.
\end{quote}

\subsection*{You'll want to know: \sk{विलम्बयति}}
\textbf{\sk{विलम्बयति}} — \emph{vilambayati} — \textquotedblleft{}delays (something); holds up\textquotedblright{}.

\begin{quote}
\sk{यदि} \sk{वर्षा} \sk{कार्यं} \sk{विलम्बयति}, \sk{तर्हि} \sk{अहं} \sk{श्वः} \sk{करिष्यामि।}
\end{quote}

\begin{quote}
\emph{yadi varṣā kāryaṃ vilambayati, tarhi ahaṃ śvaḥ kariṣyāmi.}
\end{quote}

\begin{quote}
If the rain holds up the work, I will finish tomorrow.
\end{quote}

Past (it …): \sk{वर्षा} \sk{यात्रां} \sk{विलम्बितवती}, \emph{varṣā yātrāṃ vilambitavatī}.

\begin{cousinweb}
\textbf{\sk{विलम्बयति}} (\emph{vilambayati}) is built on \emph{vi} and the root √lamb, \textquotedblleft{}to hang\textquotedblright{}, in its causative form; MW gives it the senses \textquotedblleft{}to cause to linger, to detain, to delay\textquotedblright{} (MW).
\end{cousinweb}

\begin{grammarlens}[title={What you've built}]
One more word for saying what happens if.
\end{grammarlens}

\subsection*{Your turn}
\begin{itemize}
\raggedright
  \item \emph{Say it:} \emph{vilambayati}
  \item \emph{Say it:} \emph{vilambayati}, once more
  \item \emph{Say it:} \emph{yadi varṣā kāryaṃ vilambayati, tarhi ahaṃ śvaḥ kariṣyāmi.} — If the rain holds up the work, I will finish tomorrow.
  \item \emph{Recall:} say the Sanskrit for a hindrance, then the Sanskrit for hinders, then say \emph{vilambayati} again
\end{itemize}

\begin{tcolorbox}[breakable,colback=teal!4,colframe=teal!35!black,title={Before you move on}]
What does \sk{विलम्बयति} mean? (\textquotedblleft{}Delays (something); holds up\textquotedblright{}.) Say it once more.
\end{tcolorbox}
\section[\texorpdfstring{\sk{अपनयति} (apanayati)}{अपनयति (apanayati)}]{\sk{अपनयति} (apanayati) — removes; takes away}

\label{lesson:SA-C372-apanayati}



\begin{quote}
Before the new one: say the Sanskrit for hinders, then the Sanskrit for holds up.
\end{quote}

\subsection*{You'll want to know: \sk{अपनयति}}
\textbf{\sk{अपनयति}} — \emph{apanayati} — \textquotedblleft{}removes; takes away\textquotedblright{}.

\begin{quote}
\sk{यदि} \sk{मार्गे} \sk{पाषाणः} \sk{अस्ति}, \sk{तर्हि} \sk{तम्} \sk{अपनयतु।}
\end{quote}

\begin{quote}
\emph{yadi mārge pāṣāṇaḥ asti, tarhi tam apanayatu.}
\end{quote}

\begin{quote}
If there is a stone on the road, remove it.
\end{quote}

Past (he …): \sk{सः} \sk{पाषाणम्} \sk{अपनीतवान्}, \emph{saḥ pāṣāṇam apanītavān}.

\begin{cousinweb}
\textbf{\sk{अपनयति}} (\emph{apanayati}) is built on \emph{apa}, \textquotedblleft{}away\textquotedblright{}, and the root √nī, \textquotedblleft{}to lead\textquotedblright{}; MW gives \emph{apa-nīta} \textquotedblleft{}led away from\textquotedblright{} and the verb the sense \textquotedblleft{}to take away, to remove\textquotedblright{} (MW).
\end{cousinweb}

\begin{grammarlens}[title={What you've built}]
One more word for saying what happens if.
\end{grammarlens}

\subsection*{Your turn}
\begin{itemize}
\raggedright
  \item \emph{Say it:} \emph{apanayati}
  \item \emph{Say it:} \emph{apanayati}, once more
  \item \emph{Say it:} \emph{yadi mārge pāṣāṇaḥ asti, tarhi tam apanayatu.} — If there is a stone on the road, remove it.
  \item \emph{Recall:} say the Sanskrit for hinders, then the Sanskrit for holds up, then say \emph{apanayati} again
\end{itemize}

\begin{tcolorbox}[breakable,colback=teal!4,colframe=teal!35!black,title={Before you move on}]
What does \sk{अपनयति} mean? (\textquotedblleft{}Removes; takes away\textquotedblright{}.) Say it once more.
\end{tcolorbox}
